\section*{Bab \ref{ch:b3:palaeomagnetism} --- Geomagnetisme dan Paleomagnetisme}

\begin{solution}{exo:b3:palaeomagnetism:1}
$\tan I = 2\tan\ang{30} = 1.1547$, sehingga $I = \ang{49.1}$. Sebaliknya, $\tan\lambda
= \tfrac12\tan(\ang{-20}) = -0.1820$, sehingga $\lambda = \ang{-10.3}$, yaitu
\ang{10.3}\,LS.
\end{solution}

\begin{solution}{exo:b3:palaeomagnetism:2}\begingroup\emergencystretch=3em\relax 
$H = \sqrt{20000^2 + 2000^2} = \qty{20100}{nT}$;
$F = \sqrt{20000^2 + 2000^2 + 40000^2} = \qty{44770}{nT}$. Deklinasi
magnetik adalah $\arctan(-2000/20000) = \ang{-5.7}$, \ang{5.7} ke barat dari utara;
inklinasinya $\arctan(40000/20100) = \ang{63.3}$, ke bawah.
\par\endgroup\end{solution}

\begin{solution}{exo:b3:palaeomagnetism:3}
$B_0 = -g_1^0 = \qty{29350}{nT}$ di khatulistiwa, dan $2B_0 = \qty{58700}{nT}$ di
kutub. $m = 4\pi a^3B_0/\mu_0 = a^3B_0/\num{e-7} = \num{2.5862e20}\times
\num{2.935e-5}/\num{e-7} = \qty{7.59e22}{A.m^2}$, sedikit lebih kecil daripada
\qty{7.69e22}{A.m^2} pada dipol penuh, yang juga memuat $g_1^1$ dan $h_1^1$.
\end{solution}

\begin{solution}{exo:b3:palaeomagnetism:4}
(a) \omterm{def:b3:palaeomagnetism:nrm}{Magnetisasi termoremanen}, yang merekam medan saat lava mendingin
melewati \omterm{def:b3:palaeomagnetism:blocking}{suhu pemblokiran} butir-butirnya. (b) \omterm{def:b3:palaeomagnetism:nrm}{Magnetisasi remanen detrital},
yang merekam medan selama lanau mengendap dan
terkunci oleh pemadatan. (c) \omterm{def:b3:palaeomagnetism:nrm}{Magnetisasi remanen kimia}, yang merekam medan
selama hematit tumbuh, saat penguburan awal dan karena itu sedikit
setelah pengendapan. (d) Remanensi viskos searah medan sekarang, yang ditambahkan pada
TRM asli oleh butir-butir dengan waktu relaksasi terpendek; ia dihilangkan oleh
langkah-langkah pertama demagnetisasi.
\end{solution}

\begin{solution}{exo:b3:palaeomagnetism:5}
Vektor satuan $(\cos I\cos D,\ \cos I\sin D,\ \sin I)$ adalah
\begin{gather*}
  (0.6545, 0.1391, 0.7431),\quad (0.6670, 0.2695, 0.6947),\\
  (0.5235, 0.1501, 0.8387),\quad (0.5985, 0.0629, 0.7986).
\end{gather*}\begingroup\emergencystretch=3em\relax 
Jumlahnya adalah $(2.4435, 0.6216, 3.0751)$, dengan panjang $R = 3.9766$. Rata-rata:
$D = \arctan(0.6216/2.4435) = \ang{14.3}$, $I = \arcsin(3.0751/3.9766) =
\ang{50.7}$. $k = 3/(4 - 3.9766) = 128$. Dengan $20^{1/3} - 1 = 1.7144$,
$\cos\alpha_{95} = 1 - (0.0234/3.9766)\times1.7144 = 0.98991$, sehingga
$\alpha_{95} = \ang{8.1}$: dengan empat spesimen, kerucut tetap lebar,
betapapun rapat kelompok itu.
\par\endgroup\end{solution}

\begin{solution}{exo:b3:palaeomagnetism:6}\begingroup\emergencystretch=3em\relax 
$\lambda = \arctan(\tfrac12\tan\ang{50}) = \ang{30.79}$, sehingga $p = \ang{59.21}$.
$\sin\lambda_p = \sin\ang{40}\cos\ang{59.21} + \cos\ang{40}\sin\ang{59.21}
\cos\ang{20} = 0.9474$, $\lambda_p = \ang{71.34}$\,LU. $\sin\beta =
\sin\ang{59.21}\sin\ang{20}/\cos\ang{71.34} = 0.9184$, $\beta = \ang{66.66}$.
Karena $\cos p = 0.5119 < \sin\ang{40}\times0.9474 = 0.6090$, kasus kedua
berlaku: $\varphi_p = 10 + 180 - 66.66 = \ang{123.34}$\,BT. Kutubnya berada di
\ang{71.3}\,LU, \ang{123.3}\,BT.
\par\endgroup\end{solution}

\begin{solution}{exo:b3:palaeomagnetism:7}
(a) $(7.69 - 8.32)/8.32 = \qty{-7.6}{\percent}$ dalam 125 tahun, atau
\qty{-6.1}{\percent} per abad. (b) $29350/12.6 = 2329$, sekitar 2300 tahun.
Ini bukan ramalan: laju itu sendiri berubah dari satu model lima tahunan ke
berikutnya, momen pernah lebih tinggi dan lebih rendah di masa lalu tanpa
membalik, dan medan membalik pada interval tak beraturan yang tidak diprediksi oleh
laju mana pun.
\end{solution}

\begin{solution}{exo:b3:palaeomagnetism:8}
Antara 20 dan \qty{300}{\celsius} vektor kehilangan $(6.4, 0.0, 13.6)$: komponen
sebesar \qty{15.0}{mA/m}, dengan $D = \ang{0.0}$ dan
$I = \arctan(13.6/6.4) = \ang{64.8}$. Di atas \qty{300}{\celsius} ia kehilangan $(25.8, 18.1, 24.6)$:
\qty{40.0}{mA/m}, $D = \arctan(18.1/25.8) = \ang{35.1}$, $I =
\arctan(24.6/31.5) = \ang{38.0}$. Yang kedua, yang menuju titik asal, adalah
ChRM; yang pertama, searah arah utara yang curam, dihilangkan pada suhu rendah,
adalah tumpang-tindih viskos searah medan masa kini.
\end{solution}

\begin{solution}{exo:b3:palaeomagnetism:9}
$\lambda = \arctan(\tfrac12\tan(\ang{-35})) = \ang{-19.3}$. Untuk $I = \ang{-40}$
dan \ang{-30}: \ang{-22.8} dan \ang{-16.1}, selang sebesar \ang{6.7}.
$p = \ang{90} - \lambda = \ang{109.3}$, $\cos^2p = 0.1096$:
$\mathrm dp = 2.5\times(1 + 0.329) = \ang{3.3}$ dan $\mathrm dm =
5\times\sin\ang{109.3}/\cos\ang{35} = \ang{5.8}$.
\end{solution}

\begin{solution}{exo:b3:palaeomagnetism:10}
(a) $\tan\lambda = \tfrac12$, $\lambda = \ang{26.6}$. $F$ adalah $B_0$ di
khatulistiwa dan $2B_0$ di kutub: rasio 2. (b) $\tan I = 2\tan\lambda$
memberikan $(1 + \tan^2I)\,\mathrm dI = 2(1 + \tan^2\lambda)\,\mathrm d\lambda$, sehingga
\[
  \frac{\mathrm dI}{\mathrm d\lambda} = \frac{2/\cos^2\lambda}{1 + 4\tan^2\lambda}
  = \frac{2}{\cos^2\lambda + 4\sin^2\lambda} = \frac{2}{1 + 3\sin^2\lambda}.
\]
Galat $\delta I$ memberikan $\delta\lambda = \tfrac12(1 + 3\sin^2\lambda)\,\delta I$,
dari $\tfrac12\delta I$ di khatulistiwa hingga $2\,\delta I$ di kutub: galat
terbesar terdapat di dekat kutub. (c) Di dekat khatulistiwa inklinasi berubah
dua kali lebih cepat daripada lintang, sehingga galat inklinasi tertentu hanya menimbulkan setengah
galat lintang; di dekat kutub inklinasi hampir tidak berubah terhadap lintang.
\end{solution}

\begin{solution}{exo:b3:palaeomagnetism:11}
(a) $\cos p = \sin(\ang{-35})\sin\ang{70} + \cos\ang{35}\cos\ang{70}\cos\ang{150}
= -0.5390 - 0.2426 = -0.7816$, $p = \ang{141.41}$. Lalu $\tan I_x = 2\cot p =
-2.506$, $I_x = \ang{-68.2}$, dan $\lambda_x = \ang{90} - p = \ang{-51.4}$.
$\cos D_x = (\sin\ang{70} - \sin(\ang{-35})\cos p)/(\cos\ang{35}\sin p) =
0.9617$, dan karena kutub berada di timur meridian situs ($\sin\ang{150} >
0$), $D_x = \ang{15.9}$. (b) $\lambda_o = \arctan(\tfrac12\tan(\ang{-60})) =
\ang{-40.9}$. $R = 30.0 - 15.9 = \ang{14.1}$: rotasi searah jarum jam dari blok
terhadap sumbu vertikal. $P = \lambda_x - \lambda_o = -51.4 + 40.9 =
\ang{-10.5}$: formasi itu terletak \ang{10.5} lebih jauh ke utara daripada yang
diperkirakan lempeng stabil, sehingga blok itu sejak itu telah bergerak sekitar \ang{10.5},
kira-kira \qty{1170}{km} ke selatan relatif terhadap lempeng.
\end{solution}

\begin{solution}{exo:b3:palaeomagnetism:12}
Dengan $\mathrm d\Omega = 2\pi\sin\theta'\,\mathrm d\theta'$,
\[
  \int_0^\theta\frac{\kappa e^{\kappa\cos\theta'}}{4\pi\sinh\kappa}\,2\pi\sin\theta'\,
  \mathrm d\theta' = \frac{e^\kappa - e^{\kappa\cos\theta}}{2\sinh\kappa}
  = \frac{e^\kappa - e^{\kappa\cos\theta}}{e^\kappa - e^{-\kappa}} .
\]
(b) Untuk $\kappa$ besar, $e^{-\kappa}$ dapat diabaikan; membagi dengan $e^\kappa$
menyisakan $1 - e^{-\kappa(1 - \cos\theta)} \approx 1 - e^{-\kappa\theta^2/2}$. Nilai ini
sama dengan $1 - e^{-1} = \qty{63}{\percent}$ ketika $\kappa\theta^2/2 = 1$, $\theta =
\sqrt{2/\kappa}$ radian, atau $\ang{81}/\sqrt\kappa$. (c) $81/\sqrt{50} =
\ang{11.5}$; rumus eksak memberikan \qty{63.1}{\percent}.
\end{solution}

\begin{solution}{pb:b3:palaeomagnetism:1}
\textbf{1.} Hematit pada lapisan merah dapat tumbuh dari oksida dan hidroksida besi
jauh setelah sedimennya diendapkan, sehingga menghasilkan \omterm{def:b3:palaeomagnetism:nrm}{magnetisasi remanen kimia} dengan umur yang lebih muda itu.
Uji lipatan positif, jika lapisan itu terlipat segera setelah
pengendapan, atau uji pembalikan yang positif pada lapisan-lapisan berjarak rapat dari kedua
polaritas, akan menunjukkan bahwa remanen itu terbentuk lebih awal.

\textbf{2.} $(\cos\ang{37}\cos\ang{50}, \cos\ang{37}\sin\ang{50}, \sin\ang{37})
= (0.5134, 0.6118, 0.6018)$.

\textbf{3.} $R = \sqrt{3.8943^2 + 2.6500^2 + 3.5717^2} = 5.9114$. $D =
\arctan(2.6500/3.8943) = \ang{34.2}$, $I = \arcsin(3.5717/5.9114) = \ang{37.2}$.

\textbf{4.} $k = 5/(6 - 5.9114) = 56.5$. Dengan $20^{1/5} - 1 = 0.8206$,
$\cos\alpha_{95} = 1 - (0.0886/5.9114)\times0.8206 = 0.98771$, $\alpha_{95} =
\ang{9.0}$.

\textbf{5.} B: $k = 6/(7 - 6.8979) = 58.8$; $20^{1/6} - 1 = 0.6475$,
$\cos\alpha_{95} = 1 - (0.1021/6.8979)\times0.6475 = 0.99042$, $\alpha_{95} =
\ang{7.9}$. C: $k = 7/(8 - 7.8707) = 54.1$; $20^{1/7} - 1 = 0.5340$,
$\cos\alpha_{95} = 1 - (0.1293/7.8707)\times0.5340 = 0.99123$, $\alpha_{95} =
\ang{7.6}$.

\textbf{6.} B memiliki $k$ terbesar, kelompok yang paling rapat; C memiliki
$\alpha_{95}$ terkecil, rata-rata yang paling terdefinisi, karena memiliki situs terbanyak. Ketiganya
berdekatan: ketiga formasi itu terkelompok dengan baik.

\textbf{7.} $\lambda_A = \arctan(\tfrac12\tan\ang{37.17}) = \ang{20.76}$\,LU,
$\lambda_B = \ang{32.15}$\,LU, $\lambda_C = \ang{45.33}$\,LU.

\textbf{8.} $p = 90 - 20.76 = \ang{69.24}$. $\sin\lambda_p =
\sin\ang{47}\cos\ang{69.24} + \cos\ang{47}\sin\ang{69.24}\cos\ang{34.23} =
0.2593 + 0.5272 = 0.7865$: $\lambda_p = \ang{51.86}$\,LU.

\textbf{9.} $\sin\beta = \sin\ang{69.24}\sin\ang{34.23}/\cos\ang{51.86} =
0.8516$, $\beta = \ang{58.40}$. Karena $\cos p = 0.3545 < \sin\ang{47}\times
0.7865 = 0.5752$, kasus kedua berlaku, $\cos\beta < 0$:
$\varphi_p = 15.00 + 180 - 58.40 = \ang{136.60}$\,BT.

\textbf{10.} $\cos^2p = 0.1257$: $\mathrm dp = 4.50\times1.377 = \ang{6.2}$;
$\mathrm dm = 8.99\times\sin\ang{69.24}/\cos\ang{37.17} = \ang{10.5}$.

\textbf{11.} Dari \ang{51.9}\,LU pada \qty{60}{Ma} hingga \ang{66.5}\,LU pada
\qty{30}{Ma} dan \ang{87.1}\,LU pada \qty{4}{Ma}, di meridian yang hampir sama
(\ang{136}--\ang{139}\,BT): lintasan menuju kutub geografis, dicapai dalam
kerucutnya oleh formasi termuda.

\textbf{12.} Lapisan abu, pada \qty{22}{m}, terletak di zona normal kedua,
19.7--\qty{24.6}{m}; dengan umur sekitar \qty{3.2}{Ma}, zona itu adalah C2An.2n
(3.207--3.116). Selanjutnya zona 4.2--\qty{14.8}{m} adalah C2An.3n, zona
29.2--\qty{54.0}{m} adalah C2An.1n, dasar yang terbalik termasuk kron di bawah
C2An, dan puncak yang terbalik termasuk kron di atasnya.

\textbf{13.} $(19.7 - 4.2)/(3.596 - 3.207) = 15.5/0.389 = \qty{39.8}{m/Ma}$;
$(54.0 - 24.6)/(3.116 - 2.581) = 29.4/0.535 = \qty{55.0}{m/Ma}$.

\textbf{14.} Dasar: $3.596 + 4.2/39.8 = \qty{3.70}{Ma}$. Puncak: $2.581 -
10.0/55.0 = \qty{2.40}{Ma}$.

\textbf{15.} Hanya batas pembalikan yang ditentukan penanggalannya, terpisah beberapa ratus
ribu tahun; di antara keduanya laju diasumsikan konstan. Perubahan penurunan,
pasokan sedimen atau iklim, atau celah dalam pengendapan, akan mengubahnya, dan sebuah
hiatus bahkan dapat menghilangkan zona pendek dan menyesatkan pencocokan.

\textbf{16.} $\cos p = \sin\ang{47}\sin\ang{78} + \cos\ang{47}\cos\ang{78}
\cos\ang{175} = 0.7154 - 0.1413 = 0.5741$, $p = \ang{54.96}$. $\tan I_x =
2\cot p = 1.4023$, $I_x = \ang{54.5}$; $\lambda_x = \ang{35.04}$\,LU.

\textbf{17.} $\cos D_x = (\sin\ang{78} - \sin\ang{47}\times0.5741)/
(\cos\ang{47}\sin\ang{54.96}) = 0.99976$; kutub terletak di timur
meridian situs ($\sin\ang{175} > 0$), sehingga $D_x = \ang{1.3}$.

\textbf{18.} \omterm{def:b3:palaeomagnetism:vgp}{Paleolintang} yang diharapkan dari $I_x$ = \ang{54.5}, \ang{60.0}, dan
\ang{63.6}: \ang{35.03}, \ang{40.89}, dan \ang{45.21}. $P_A = 35.03 - 20.76 =
\ang{14.3}$, $P_B = 40.89 - 32.15 = \ang{8.7}$, $P_C = 45.21 - 45.33 =
\ang{-0.1}$.

\textbf{19.} $R_A = 34.2 - 1.3 = \ang{32.9}$, $R_B = 23.9 - 2.9 = \ang{21.0}$,
$R_C = 3.4 - 2.5 = \ang{0.9}$, semuanya searah jarum jam.

\textbf{20.} $\Delta D_o = \arcsin(\sin\ang{8.99}/\cos\ang{37.17}) = \ang{11.3}$;
$\Delta D_x = \arcsin(\sin\ang{3.0}/\cos\ang{54.5}) = \ang{5.2}$; $\Delta R_A =
\sqrt{11.3^2 + 5.2^2} = \ang{12.4}$.

\textbf{21.} $R_A = \ang{32.9} \pm \ang{12.4}$ dan $R_B = \ang{21.0} \pm
\ang{13.8}$ signifikan; $R_C = \ang{0.9} \pm \ang{17.9}$ tidak.

\textbf{22.} Enam puluh juta tahun lalu Kesta terletak sekitar \ang{14} lintang
ke selatan dari posisi yang ditempatkan benua stabil, dan telah berputar sekitar \ang{33}
berlawanan arah jarum jam dari orientasinya sekarang. Pada \qty{30}{Ma} ia telah bergerak
\ang{6} dari \ang{14} itu ke utara dan berputar \ang{12} dari \ang{33} itu
searah jarum jam; pada \qty{4}{Ma} ia telah mencapai tempatnya sekarang di benua,
tanpa rotasi tersisa. Ia bergerak ke utara di sepanjang tepi benua, berputar
searah jarum jam selama bergerak, dan telah berada di tempatnya pada \qty{4}{Ma}.

\textbf{23.} $47.00 - 20.76 = \ang{26.24}$, atau $26.24\times111.19 =
\qty{2918}{km}$. Ini termasuk gerak benua stabil itu sendiri;
relatif terhadapnya, Kesta bergerak $P_A = \ang{14.3}$.

\textbf{24.} $I = \ang{28.18}$ dan \ang{46.16} memberikan $\lambda = \ang{15.00}$ dan
\ang{27.50}: setengah lebar \ang{6.25}, dan satu simpangan baku dari
$6.25/1.96 = \ang{3.19}$, atau \qty{355}{km}. Sejak formasi A, Kesta telah bergerak
ke utara sejauh
\textbf{\qty[text-series-to-math=true]{2918 \pm 355}{km}}.
\end{solution}
